% SPDX-License-Identifier: Apache-2.0
% covers: EthSlots
\datasheet{EthSlots}{The Ethernet slot registers on AXI-Lite: LiteEth's
map, with the frame buffers in main memory.}

\dsfacts{\code{lib/parts/src/ethslots.rs}}%
{\code{txhdl\_parts::ethslots::EthSlots}}{\code{//docs:examples}}

\subsection*{Function}

\code{EthSlots} is the register map of LiteX's LiteEth in hardware,
at LiteX's offsets, so that Linux's \code{litex\_liteeth}, whose
offsets are fixed, drives it unchanged (issue 1203), and Zephyr's
\code{drivers/ethernet/eth\_litex\_liteeth.c} takes them from the
header \code{//tools/regmap} writes. Bringing up Ethernet is a port of
a driver that already works rather than a design nobody has driven.

Its AXI-Lite side is one port, \code{bus}, a \code{LitePort} with
32-bit addresses and words: \code{bus\_aw}, \code{bus\_ar} and
\code{bus\_w} in, and \code{bus\_b} and \code{bus\_r} out. Ten
further ports face the engines that move the bytes and the unit that
receives them: \code{tx\_busy}, \code{rx\_busy}, \code{rx\_len},
\code{rx\_which} and \code{rx\_drops} in, and \code{tx\_base},
\code{tx\_bytes}, \code{tx\_start}, \code{rx\_base} and
\code{rx\_full} out. \code{irq} is high while a received frame is
unacknowledged and the receive interrupt is enabled.

It holds no frame's bytes. There are two slots each way, used in turn,
and each is a flat buffer of 2048 bytes in main memory rather than
inside the peripheral, so that \code{LineFetch} and \code{LineStore}
move the bytes instead of the processor. The unit's work is to say
which buffer and how many bytes, to tell the driver when a frame has
landed, and to keep the frames that have landed and not been
acknowledged in order, a slot and a length each, at most one a slot
(issue 1313).

Two slots a direction rather than a descriptor ring, because Zephyr's
driver model never hands a chain to hardware:
\code{net\_pkt\_read(pkt, buf, len)} and \code{net\_pkt\_write(pkt,
buf, len)} take a flat pointer and a byte count, so the hardware only
ever sees one contiguous region per frame.

\subsection*{Parameters}

\code{BASE}, where the four buffers begin. A slot is \code{BASE + n *
2048}, receive slots first and transmit slots 4096 bytes above them.
The board uses \code{0x4100\_0000}, which \code{isa::ETH\_BUF\_BASE}
states.

The two directions are separate regions and have to be. A slot number
alone would make transmit slot zero and receive slot zero the same
address, and a frame arriving would land on one waiting to go out.

\textbf{\code{BASE} must be 1 KB aligned, and nothing checks it.} A
slot is then 1 KB aligned too, since the slots are 2 KB apart, and a
256 beat burst of words is exactly 1 KB, so no burst can cross AXI4's
4 KB boundary. An unaligned \code{BASE} puts that obligation on the
caller, and a burst that straddles a boundary is a protocol violation
rather than a wrong address, so it will not look like a fault in this
unit.

\subsection*{Register map}

The unit decodes the word from address bits 5 to 2, so the words are
at offsets \code{0x0} to \code{0x3c} of its range. The board puts
them at \code{0x3400}, which \code{isa::ETH\_BASE} states.

\dsregs{EthSlots}

The fourteen words are LiteX's, in LiteX's order, and words 14 and
15 are named by no register. Three words exist for the drivers and
hold nothing here: \code{tx\_level} reads zero, since one transmit
runs at a time, and \code{tx\_ev\_status} reads zero, since the
transmit raises no event. \code{rx\_errors} reads \code{rx\_drops},
the frames the receiving side dropped because both slots held a frame
the driver had not acknowledged; it read zero, and said none were
dropped, before issue 1313. \code{rx\_ev\_status} reads the same bit as
\code{rx\_ev\_pending}. \code{tx\_ev\_pending} is never set, as
the behaviour below says, and \code{tx\_ev\_enable} is kept so the
driver can write it.

The write-only words, \code{tx\_start}, \code{tx\_slot} and
\code{tx\_length}, and the unnamed 14 and 15 read zero. They once
read as \code{tx\_ev\_enable}, the read path's fall-through, so a
driver that wrote \code{tx\_length} and read it back got an unrelated
bit (issue 454); every word the read names now, and the rest are
zero.

\emph{The transmit side has no interrupt.} LiteEth's driver disables
it and polls \code{tx\_ready}, because Zephyr's send may block, so
\code{tx\_ev\_pending} is never set by anything and
\code{tx\_ev\_enable} gates nothing. Both words exist because the
driver acknowledges them. A driver that waited on
\code{tx\_ev\_pending} would wait for ever.

\dstables{EthSlots}

\subsection*{Behaviour}

\begin{itemize}
\item A read is answered in the cycle the request is taken, and a
write is answered with \code{bus\_b}. The unit never stalls: it needs
nothing the other AXI-Lite peripherals do not.
\item \textbf{An arrival is the store engine falling idle, not the
frame being received.} \code{LineStore} holds \code{running} high
until the write response comes back from memory, so its fall is the
first moment the frame is certainly there. The unit takes that level
rather than a completion pulse, deliberately: a pulse could be joined
to the receiver by somebody reading the map and not the comment, and
the interrupt would then tell a driver to read a frame whose last
beats are still in flight.
\item \textbf{Frames wait in order (issue 1313).} An arrival puts
\code{rx\_which} and \code{rx\_len} behind any frame already waiting.
\code{rx\_slot} and \code{rx\_length} read the oldest, and
\code{rx\_ev\_pending} reads one while any waits. Writing one to it
takes the oldest off, so a driver that acknowledges one frame an
interrupt, as Zephyr's and Linux's do, is interrupted again for the
next. Before issue 1313 an arrival replaced the slot and length the
driver had been told of, and a driver busy with one frame was told of
the last of a burst only.
\item \textbf{An acknowledgement in the same cycle as an arrival does
not lose the frame.} The oldest frame is taken off and the new one put
behind what is left, so the second frame comes to the front.
\item \textbf{A frame with no slot is dropped and counted.}
\code{rx\_full} is high while both slots hold a frame, or while one
does and the frame just stored is about to be counted. The receiving
side, \code{FrameIn}, then takes the next frame's bytes, stores none
of them, and counts it, and \code{rx\_errors} reads the count. A
frame dropped at the input is one the sender retransmits; a frame
landed on one being read was two frames lost, silently.
\item \code{tx\_start} is a request rather than a pulse the engine
must catch. Writing one sets a flag; the unit drives \code{tx\_start}
high while the flag is set and the engine is not busy, and clears the
flag when the engine goes busy. So the request lasts until it is
taken.
\item \code{tx\_base} is \code{BASE + 4096 + tx\_slot * 2048} and
\code{rx\_base} is \code{BASE + rx\_which * 2048}. The receiving side
is told where the slot it is filling begins, and which slot that is
comes from the engine, so that a frame lands somewhere the driver is
not reading.
\item \code{irq} is high while any frame waits and
\code{rx\_ev\_enable} is set. It is a level, and the board's
interrupt controller asks again for a level still high on the
complete.
\end{itemize}

\subsection*{What a driver must do}

\textbf{Read the last word of a frame back before writing
\code{tx\_start}.} Stores on this machine are posted, so without
an order the fetch engine may read the slot before the frame has
landed and send what was there before. A \code{fence} now waits
until every posted store has been answered (issue 432), and either
the fence or the read orders them; the read is what was proven on
the board, and \code{ethtx.rs} keeps it.

What makes the obligation dischargeable is measured rather than
assumed: this link answers a load behind an unanswered store to the
same address with the stored value, which
\code{a\_load\_behind\_an\_unanswered\_store\_to\_one\_address} in
\code{lib/parts/src/bus/axi/sim.rs} asserts, and \code{//docs:axi}
states.

\textbf{That is this interconnect's behaviour and not AXI's
guarantee.} AXI does not order a read behind a write, and a driver
written for another fabric on the strength of this paragraph would be
wrong. The reason the read works here is the link, not the memory
being a serialisation point.

\subsection*{Verification}

\begin{itemize}
\item \code{only\_word\_13\_reads\_tx\_ev\_enable}, in
\code{//lib/parts:parts\_test}, sets \code{tx\_ev\_enable} and reads
every word, so a word that mirrored it again would fail (issue 454).
\item \code{ex\_ethslots} drives the map as
\code{eth\_litex\_liteeth.c} will, in the order that driver does it:
\code{tx\_ready} one while idle and zero while a frame is going, no
interrupt while the store is in flight and one after it, and the
pending bit tested both ways, since a write of zero must not
acknowledge.
\item It also tests the cycle where an acknowledgement meets an
arrival. The cycle to issue at was measured against both orders of
the drives rather than reasoned about, because two earlier attempts
at that test passed against the faulty order.
\item It lands a third frame while the second waits, and checks that
it waits behind it, that \code{rx\_full} is up, that
\code{rx\_errors} reads the receiving side's count, and that two
acknowledgements take the two frames off in order (issue 1313).
\item The build simulates the netlist against that run under nvc and
Verilator: \code{//docs:ethslots\_sim\_ethslots\_tb\_test} and
\code{//docs:ethslots\_vsim\_test}.
\item On the board, with the engines behind it, the two Ethernet tests
in \code{cpu/vreteno/tests/board.rs} run programs on the core that
use its registers as the driver does: one waits for an arrival and
reads the frame from the slot and length it names, and the other
sends two frames back to back, waiting on \code{tx\_ready} between
them.
\item \code{frames\_beyond\_the\_slots\_are\_dropped\_and\_counted\_not\_lost},
in the same file, sends eight full frames back to back to a driver
that spends five thousand loop turns on each: it is told of three,
intact, and the other five are counted. Before issue 1313 it was told
of two, both overwritten under it, and \code{rx\_errors} read zero.
\end{itemize}

\subsection*{Used by}

The board, at \code{cpu/vreteno/src/board.rs}, on the fifth AXI-Lite
slot at \code{0x3400}, with its arrival as the third source of
\code{Plic3}. Its engine-facing ports go to the fetch engine and
\code{FrameOut} on the way out, and to \code{FrameIn} and the store
engine on the way in (issue 151). The example \code{ex\_ethslots}.

\subsection*{Limits}

It moves no frames itself: the engines behind it do. Without them,
as in \code{ex\_ethslots}, a driver binds and reads and writes every
register, and no frame arrives or leaves.

Two slots a direction is queue depth and not throughput: a burst
longer than two frames, while the driver is busy, loses the rest,
counted. A ring, if
it is ever wanted, is an addition rather than a rewrite: a descriptor
states the same three values \code{tx\_slot}, \code{tx\_length} and
\code{tx\_start} already state.

The unit checks no more of the address than the word, because the
bridge in front of it sends it only its own range.
